\chapter{LASIK Eye Surgery}

\section*{Overview}
LASIK is a refractive eye procedure that reshapes the cornea with a laser to reduce dependence on glasses or contact lenses.

\section*{Why It Is Done}
It is an elective treatment for suitable adults with stable nearsightedness, farsightedness, or astigmatism after a comprehensive eye evaluation.

\section*{How It Is Performed}
A thin corneal flap is created, an excimer laser reshapes the underlying cornea, and the flap is repositioned. The procedure is usually performed with topical anesthetic.

\section*{Risks and Recovery}
Dry eye, glare, halos, under- or over-correction, infection, corneal ectasia, flap problems, and need for additional treatment are possible. Vision often improves quickly, but stabilization takes longer.

\section*{References}
\begin{enumerate}
  \item National Eye Institute. LASIK. 2025.
  \item Current Medical Diagnosis and Treatment. McGraw-Hill; 2025.
\end{enumerate}
\clearpage
